\section{OBJETIVOS}
\subsection{Objetivo general}
Analizar mediante simulación electromagnética en ANSYS HFSS el comportamiento de un combinador de guía de onda WR42 operando a 20 GHz, evaluando su desempeño a partir de los parámetros de dispersión y la distribución del campo eléctrico dentro de la estructura.
\subsection{Objetivos especificos}
\begin{itemize}
    \item Construir la geometría tridimensional del combinador de guía de onda a partir de las dimensiones establecidas para una estructura WR42.
    \item Aplicar condiciones de frontera y excitaciones adecuadas que representen el comportamiento físico del dispositivo durante la simulación.
    \item Configurar el análisis modal y el barrido de frecuencia necesario para obtener la respuesta electromagnética del combinador alrededor de la frecuencia de operación.
    \item Obtener y analizar los parámetros S del sistema, especialmente los relacionados con reflexión, transmisión, aislamiento y combinación de potencia entre los puertos.
    \item Visualizar la distribución del campo eléctrico dentro de la guía de onda para identificar las zonas de mayor concentración de energía y comprender el comportamiento de la unión combinadora.
    \item Relacionar los resultados obtenidos en la simulación con el funcionamiento esperado de un combinador de potencia de baja pérdida en aplicaciones de microondas.
\end{itemize}
